%% toc: P1 5 5.7
\section{Proofs for transfinite recursion on an ordinal}
\label{sec:annex-p01-w02-p01-013-transfinite-recursion}

Notation and readings are those of \cref{def:p01-w02-p01-013-reading}.
Throughout, \(\delta\) is an ordinal and \(G\) is a functional rule for
\(\delta\), given by the formula \(\theta(g,y,\bar p)\). Item numbers
``item~1'' to ``item~6'' without further reference are the items of
\cref{thm:p01-w02-p01-011-ordinals}. The strict well-ordering of an ordinal
\(\gamma\) by \(\in\) means, by the definition in
\cref{sec:set-prerequisites-03}, that every nonempty \(S\subseteq\gamma\) has
a member \(\alpha\) such that no \(\beta\in S\) satisfies \(\beta\in\alpha\);
such an \(\alpha\) is called an \(\in\)-least member of \(S\).

Two facts about equality are used repeatedly. First, two functions \(f\) and
\(f'\) with the same domain and with \(f(x)=f'(x)\) at every \(x\) of that
domain are equal: a member of \(f\) is a pair \((x,f(x))\) with \(x\) in the
domain, hence equals \((x,f'(x))\in f'\), and symmetrically, so
Extensionality gives \(f=f'\). Second, if \(g=g'\) are functions in the domain
of the rule, then \(G(g)=G(g')\): the unique \(y\) with \(\theta(g,y,\bar p)\)
is the unique \(y\) with \(\theta(g',y,\bar p)\), by substitutivity of
identity.

\begin{annexitem}{annex:p01-w02-p01-013-agreement}{thm:p01-w02-p01-013-agreement}{Restriction and agreement of attempts}
\begin{proof}
\emph{Item~1.} Let \(\gamma\leq\delta\). If \(\gamma=\delta\), then
\(\gamma\) is an ordinal and \(\gamma\subseteq\delta\). If
\(\gamma\in\delta\), then \(\gamma\) is an ordinal by item~1 of
\cref{thm:p01-w02-p01-011-ordinals}, and \(\gamma\subseteq\delta\) because
\(\delta\) is transitive. Let \(\alpha\in\gamma\). Then \(\alpha\) is an
ordinal by that item~1 applied to \(\gamma\); \(\alpha\in\delta\) because
\(\gamma\subseteq\delta\), that is, \(\alpha<\delta\); and
\(\alpha\subseteq\gamma\) because \(\gamma\) is transitive.

Let \(f\) be a function with domain \(\gamma\) and \(A\subseteq\gamma\). The
set \(f\mathbin{\upharpoonright}A\) is a subset of \(f\), so no two of its
pairs share a first coordinate, and it is a function. Every first coordinate
of a pair in it lies in \(A\) by definition. Conversely, for \(x\in A\) we
have \(x\in\gamma\), so the pair \((x,f(x))\) lies in \(f\), and since
\(x\in A\) it lies in \(f\mathbin{\upharpoonright}A\). Hence
\(\dom(f\mathbin{\upharpoonright}A)=A\) and
\((f\mathbin{\upharpoonright}A)(x)=f(x)\) for \(x\in A\). If
\(B\subseteq A\), a pair \(p\) belongs to
\((f\mathbin{\upharpoonright}A)\mathbin{\upharpoonright}B\) iff \(p\in f\)
and its first coordinate lies in \(A\) and in \(B\), iff \(p\in f\) and its
first coordinate lies in \(B\), because \(B\subseteq A\). Extensionality gives
\((f\mathbin{\upharpoonright}A)\mathbin{\upharpoonright}B
=f\mathbin{\upharpoonright}B\).

In particular, for an attempt \(f\) of length \(\gamma\) and
\(\alpha<\gamma\), the function \(f\mathbin{\upharpoonright}\alpha\) has
domain the ordinal \(\alpha<\delta\), so \(G(f\mathbin{\upharpoonright}\alpha)\)
is defined. This justifies \cref{def:p01-w02-p01-013-attempt}.

\emph{Item~2.} Let \(f\) and \(f'\) be attempts of lengths \(\gamma\) and
\(\gamma'\). By Separation form
\[
  D=\{\alpha\in\gamma:\alpha\in\gamma'\text{ and }f(\alpha)\ne f'(\alpha)\}.
\]
Suppose \(D\ne\varnothing\). Since \(D\subseteq\gamma\subseteq\delta\), the
well-ordering of \(\delta\) gives an \(\in\)-least member \(\alpha\) of
\(D\). Then \(\alpha\in\gamma\) and \(\alpha\in\gamma'\), so by item~1
\(\alpha\subseteq\gamma\) and \(\alpha\subseteq\gamma'\). Let
\(\beta\in\alpha\). Then \(\beta\in\gamma\) and \(\beta\in\gamma'\), and
\(\beta\notin D\) because \(\beta\in\alpha\) and \(\alpha\) is
\(\in\)-least in \(D\). Hence \(f(\beta)=f'(\beta)\). By item~1 the functions
\(f\mathbin{\upharpoonright}\alpha\) and \(f'\mathbin{\upharpoonright}\alpha\)
both have domain \(\alpha\), and at each \(\beta\in\alpha\)
\[
  (f\mathbin{\upharpoonright}\alpha)(\beta)=f(\beta)=f'(\beta)
  =(f'\mathbin{\upharpoonright}\alpha)(\beta).
\]
So \(f\mathbin{\upharpoonright}\alpha=f'\mathbin{\upharpoonright}\alpha\), and
\[
  f(\alpha)=G(f\mathbin{\upharpoonright}\alpha)
  =G(f'\mathbin{\upharpoonright}\alpha)=f'(\alpha),
\]
the outer equalities being the attempt conditions at \(\alpha\). This
contradicts \(\alpha\in D\). Hence \(D=\varnothing\): \(f(\alpha)=f'(\alpha)\)
whenever \(\alpha\in\gamma\) and \(\alpha\in\gamma'\).

\emph{Item~3.} Two attempts \(f\) and \(f'\) of the same length \(\gamma\)
have the same domain \(\gamma\) and, by item~2, the same value at every
\(\alpha\in\gamma\). So \(f=f'\).

\emph{Item~4.} Let \(f\) be an attempt of length \(\gamma\) and
\(\beta\leq\gamma\). If \(\beta=\gamma\), then \(\beta\subseteq\gamma\); if
\(\beta\in\gamma\), then \(\beta\subseteq\gamma\) by item~1. Also
\(\beta\leq\delta\): either \(\beta=\gamma\leq\delta\), or
\(\beta\in\gamma\subseteq\delta\). By item~1,
\(f\mathbin{\upharpoonright}\beta\) is a function with domain \(\beta\). Let
\(\alpha\in\beta\). Then \(\alpha\in\gamma\), and \(\alpha\subseteq\beta\) by
item~1 applied to \(\beta\leq\delta\). Hence, using item~1 twice,
\[
  (f\mathbin{\upharpoonright}\beta)(\alpha)=f(\alpha)
  =G(f\mathbin{\upharpoonright}\alpha)
  =G\bigl((f\mathbin{\upharpoonright}\beta)\mathbin{\upharpoonright}\alpha\bigr).
\]
So \(f\mathbin{\upharpoonright}\beta\) is an attempt of length \(\beta\).
\end{proof}
\end{annexitem}

\begin{annexitem}{annex:p01-w02-p01-013-assembly}{thm:p01-w02-p01-013-assembly}{Assembling the earlier attempts}
\begin{proof}
Let \(\gamma\leq\delta\), and suppose that for every \(\alpha\in\gamma\) an
attempt of length \(\alpha\) exists. Each such \(\alpha\) satisfies
\(\alpha\leq\delta\) by \cref{thm:p01-w02-p01-013-agreement}, item~1, so
attempts of length \(\alpha\) are defined, and by item~3 of that lemma there
is exactly one; call it \(f_\alpha\).

\emph{The set \(h_\gamma\).} The property ``\(f\) is an attempt of length
\(\alpha\)'' is expressed by a formula of the set-theoretic language with
parameters \(\delta,\bar p\): it says that \(f\) is a function, that its
first coordinates are exactly the members of \(\alpha\), and that
\(\theta(f\mathbin{\upharpoonright}\beta,f(\beta),\bar p)\) holds for every
\(\beta\in\alpha\). Let \(\psi(\alpha,z)\) be the formula
\[
  \exists f\,\exists y\,\bigl(f\text{ is an attempt of length }\alpha
  \land\theta(f,y,\bar p)\land z=(\alpha,y)\bigr).
\]
For each \(\alpha\in\gamma\) there is exactly one \(z\) with
\(\psi(\alpha,z)\). There is at least one: \(f_\alpha\) is an attempt of
length \(\alpha\), its domain is an ordinal smaller than \(\delta\), so there
is a \(y\) with \(\theta(f_\alpha,y,\bar p)\), namely \(y=G(f_\alpha)\), and
\(z=(\alpha,y)\) exists by Pairing. There is at most one: if
\(\psi(\alpha,z)\) is witnessed by \(f,y\) and \(\psi(\alpha,z')\) by
\(f',y'\), then \(f=f'=f_\alpha\) by item~3 of
\cref{thm:p01-w02-p01-013-agreement}, so \(y=y'\) because \(\theta\) gives
exactly one value at \(f_\alpha\), and hence \(z=z'\).

The Replacement schema, applied to the set \(\gamma\) and the formula
\(\psi\), gives a set \(h_\gamma\) whose members are exactly the \(z\) with
\(\psi(\alpha,z)\) for some \(\alpha\in\gamma\), that is, the pairs
\((\alpha,G(f_\alpha))\) with \(\alpha\in\gamma\).

\emph{\(h_\gamma\) is a function with domain \(\gamma\).} Every member of
\(h_\gamma\) is an ordered pair whose first coordinate lies in \(\gamma\),
and every \(\alpha\in\gamma\) is the first coordinate of the member
\((\alpha,G(f_\alpha))\). If \((\alpha,y)\) and \((\alpha,y')\) are members,
then \((\alpha,y)=(\alpha_1,G(f_{\alpha_1}))\) and
\((\alpha,y')=(\alpha_2,G(f_{\alpha_2}))\) for some
\(\alpha_1,\alpha_2\in\gamma\); equality of ordered pairs gives
\(\alpha_1=\alpha=\alpha_2\), and then \(y=G(f_\alpha)=y'\). So \(h_\gamma\)
is a function, \(\dom h_\gamma=\gamma\), and \(h_\gamma(\alpha)=G(f_\alpha)\)
for \(\alpha\in\gamma\).

\emph{Restrictions.} Let \(\alpha\in\gamma\). By item~1 of
\cref{thm:p01-w02-p01-013-agreement}, \(\alpha\subseteq\gamma\), and the
functions \(h_\gamma\mathbin{\upharpoonright}\alpha\) and \(f_\alpha\) both
have domain \(\alpha\). Let \(\beta\in\alpha\). Then \(\beta\in\gamma\),
because \(\gamma\) is transitive, so \(f_\beta\) is defined. By item~4 of
that lemma, \(f_\alpha\mathbin{\upharpoonright}\beta\) is an attempt of
length \(\beta\), so \(f_\alpha\mathbin{\upharpoonright}\beta=f_\beta\) by
item~3. Therefore
\[
  f_\alpha(\beta)=G(f_\alpha\mathbin{\upharpoonright}\beta)=G(f_\beta)
  =h_\gamma(\beta)=(h_\gamma\mathbin{\upharpoonright}\alpha)(\beta),
\]
using the attempt condition for \(f_\alpha\) at \(\beta\), then the
definition of \(h_\gamma\), then item~1. Two functions with domain
\(\alpha\) that agree at every point are equal, so
\(h_\gamma\mathbin{\upharpoonright}\alpha=f_\alpha\).

\emph{\(h_\gamma\) is an attempt of length \(\gamma\).} For
\(\alpha\in\gamma\),
\[
  h_\gamma(\alpha)=G(f_\alpha)=G(h_\gamma\mathbin{\upharpoonright}\alpha).
\]
If \(\gamma=0\), the hypothesis is vacuous, \(h_\gamma=\varnothing\), and
\(\varnothing\) is the attempt of length \(0\).
\end{proof}
\end{annexitem}

\begin{annexitem}{annex:p01-w02-p01-013-recursion}{thm:p01-w02-p01-013-recursion}{Transfinite recursion}
\begin{proof}
An attempt of length \(\delta\) is, by
\cref{def:p01-w02-p01-013-attempt}, a function \(F\) with domain \(\delta\)
such that \(F(\alpha)=G(F\mathbin{\upharpoonright}\alpha)\) for every
\(\alpha<\delta\). So the theorem says that there is exactly one attempt of
length \(\delta\).

\emph{Attempts of every smaller length.} By Separation form
\[
  C=\{\gamma\in\delta:\text{there is no attempt of length }\gamma\};
\]
the condition is a formula with parameters \(\delta,\bar p\), as noted in
the proof of \cref{thm:p01-w02-p01-013-assembly}. Suppose
\(C\ne\varnothing\), and let \(\gamma\) be an \(\in\)-least member of \(C\)
in the well-ordering of \(\delta\). Let \(\alpha\in\gamma\). Then
\(\alpha\in\delta\), because \(\delta\) is transitive, and
\(\alpha\notin C\), because \(\alpha\in\gamma\) and \(\gamma\) is
\(\in\)-least in \(C\). So there is an attempt of length \(\alpha\). Since
\(\gamma\in\delta\), \(\gamma\leq\delta\), and
\cref{thm:p01-w02-p01-013-assembly} yields an attempt \(h_\gamma\) of length
\(\gamma\). This contradicts \(\gamma\in C\). Hence \(C=\varnothing\): for
every \(\gamma\in\delta\) there is an attempt of length \(\gamma\).

\emph{Existence.} Apply \cref{thm:p01-w02-p01-013-assembly} with
\(\gamma=\delta\). Its hypothesis was just established, so
\(F=h_\delta\) is an attempt of length \(\delta\).

\emph{Uniqueness.} By item~3 of \cref{thm:p01-w02-p01-013-agreement} with
\(\gamma=\delta\), there is at most one attempt of length \(\delta\).

\emph{What was used.} Separation formed \(D\), \(C\) and the restrictions;
Pairing formed the ordered pairs; Replacement formed \(h_\gamma\);
Extensionality gave the equalities of functions. Every object introduced is
the unique object with a stated property, or an arbitrary member of a set
already given, so no choice principle is used. The case \(\delta=0\) is
included: then \(C=\varnothing\) and \(F=\varnothing\).
\end{proof}
\end{annexitem}

\begin{annexitem}{annex:p01-w02-p01-013-coherence}{thm:p01-w02-p01-013-coherence}{Shorter lengths give restrictions}
\begin{proof}
Let \(\gamma\leq\delta\). A function whose domain is an ordinal
\(\alpha<\gamma\) has \(\alpha<\delta\) by
\cref{thm:p01-w02-p01-013-agreement}, item~1, so \(G\) assigns it a set. Thus
\(G\), given by the same formula \(\theta\), is a functional rule for
\(\gamma\), and \cref{thm:p01-w02-p01-013-recursion} for \(\gamma\) yields a
unique \(F_\gamma\) with domain \(\gamma\) and
\(F_\gamma(\alpha)=G(F_\gamma\mathbin{\upharpoonright}\alpha)\) for
\(\alpha<\gamma\).

The function \(F\) is an attempt of length \(\delta\), so by item~4 of
\cref{thm:p01-w02-p01-013-agreement} the function
\(F\mathbin{\upharpoonright}\gamma\) is an attempt of length \(\gamma\): its
domain is \(\gamma\) and
\((F\mathbin{\upharpoonright}\gamma)(\alpha)
=G((F\mathbin{\upharpoonright}\gamma)\mathbin{\upharpoonright}\alpha)\) for
every \(\alpha<\gamma\). This is the condition that defines \(F_\gamma\), and
the uniqueness clause of the theorem for \(\gamma\) gives
\(F_\gamma=F\mathbin{\upharpoonright}\gamma\).
\end{proof}
\end{annexitem}

\begin{annexitem}{annex:p01-w02-p01-013-cases}{thm:p01-w02-p01-013-cases}{Recursion by zero, successor and limit clauses}
\begin{proof}
\emph{The three kinds of ordinal.} By the definitions following
\cref{def:ordinal-cardinal}, a successor ordinal is one of the form
\(\eta\cup\{\eta\}=\eta+1\) for an ordinal \(\eta\), and a limit ordinal is a
nonzero ordinal that is not a successor. The ordinal \(0=\varnothing\) is not
a successor, because \(\eta\in\eta+1\) makes every successor nonempty. Hence
every ordinal is exactly one of: \(0\); a successor; a limit. If \(\alpha\)
is a successor, the \(\eta\) with \(\alpha=\eta+1\) is unique: if
\(\eta+1=\eta'+1\) with \(\eta\ne\eta'\), then \(\eta\in\eta'+1\) and
\(\eta\ne\eta'\) give \(\eta\in\eta'\), and symmetrically
\(\eta'\in\eta\), which item~3 of \cref{thm:p01-w02-p01-011-ordinals}
excludes. In that case \(\eta\in\alpha\), so \(\eta<\alpha\).

\emph{The combined rule.} For a function \(g\) whose domain is an ordinal
\(\alpha<\delta\), put
\[
  G(g)=
  \begin{cases}
    a & \text{if }\alpha=0,\\
    S(\eta,g(\eta)) & \text{if }\alpha=\eta+1,\\
    L(g) & \text{if }\alpha\text{ is a limit ordinal.}
  \end{cases}
\]
Written out, \(\theta(g,y)\) is the disjunction of three formulas, one for
each case, each of which states the case condition together with the
corresponding equation for \(y\); the parameters are \(a\) and those of
\(S\) and \(L\). For each such \(g\) exactly one case condition holds, and
in that case exactly one \(y\) satisfies the equation: in the first case
\(y=a\); in the second, \(\eta\) is unique, \(\eta\in\alpha=\dom g\) so
\(g(\eta)\) is defined, and \(\eta<\delta\) because \(\eta\in\alpha\) and
\(\alpha\subseteq\delta\), so \(S(\eta,g(\eta))\) is defined; in the third,
\(L(g)\) is defined because \(\dom g\) is a limit ordinal smaller than
\(\delta\). So \(G\) is a functional rule for \(\delta\).

\emph{Existence.} Let \(F\) be the function given by
\cref{thm:p01-w02-p01-013-recursion} for \(\delta\) and \(G\). If
\(0<\delta\), the function \(F\mathbin{\upharpoonright}0\) has domain \(0\),
so \(F(0)=G(F\mathbin{\upharpoonright}0)=a\). If \(\eta+1<\delta\), the
function \(F\mathbin{\upharpoonright}(\eta+1)\) has domain \(\eta+1\), and
the second case gives
\[
  F(\eta+1)=S\bigl(\eta,(F\mathbin{\upharpoonright}(\eta+1))(\eta)\bigr)
  =S(\eta,F(\eta)),
\]
using \(\eta\in\eta+1\) and item~1 of
\cref{thm:p01-w02-p01-013-agreement}. If \(\lambda<\delta\) is a limit
ordinal, the third case gives
\(F(\lambda)=L(F\mathbin{\upharpoonright}\lambda)\).

\emph{Uniqueness.} Let \(F'\) be a function with domain \(\delta\)
satisfying the three clauses, and let \(\alpha<\delta\). If \(\alpha=0\),
then \(G(F'\mathbin{\upharpoonright}0)=a=F'(0)\). If \(\alpha=\eta+1\), then
\(G(F'\mathbin{\upharpoonright}\alpha)
=S(\eta,(F'\mathbin{\upharpoonright}\alpha)(\eta))=S(\eta,F'(\eta))
=F'(\alpha)\). If \(\alpha\) is a limit, then
\(G(F'\mathbin{\upharpoonright}\alpha)=L(F'\mathbin{\upharpoonright}\alpha)
=F'(\alpha)\). So \(F'(\alpha)=G(F'\mathbin{\upharpoonright}\alpha)\) for
every \(\alpha<\delta\), and the uniqueness clause of
\cref{thm:p01-w02-p01-013-recursion} gives \(F'=F\).
\end{proof}
\end{annexitem}

\begin{annexitem}{annex:p01-w02-p01-013-closure-stages}{thm:p01-w02-p01-013-closure-stages}{Values in the closure example}
\begin{proof}
Statements about all \(n\in\N\) are proved by induction on \(n\), using that
\(\N\) is the least set containing \(0\) and closed under \(n\mapsto n+1\)
(\cref{sec:set-prerequisites-02}). By item~6 of
\cref{thm:p01-w02-p01-011-ordinals}, every \(n\in\N\) is an ordinal.

\emph{Every \(n\in\N\) is \(0\) or \(m+1\) with \(m\in\N\).} The set of
\(n\in\N\) with this property contains \(0\) and contains \(n+1\) for each of
its members \(n\), with \(m=n\); so it is \(\N\).

\emph{\(\N\) is a limit ordinal.} The set of \(n\in\N\) with
\(n\subseteq\N\) contains \(0=\varnothing\), and if \(n\subseteq\N\) then
\(n+1=n\cup\{n\}\subseteq\N\) because \(n\in\N\); so every \(n\in\N\) is a
subset of \(\N\), and \(\N\) is transitive. It is a transitive set of
ordinals, hence an ordinal by item~5. It is nonzero because \(0\in\N\). If
\(\N=\eta+1\), then \(\eta\in\N\), so \(\eta+1\in\N\), that is
\(\N\in\N\), against item~1. So \(\N\) is a limit ordinal, and \(\N+1\) and
\((\N+1)+1\) are ordinals by item~5. Put \(\delta=(\N+1)+1\); its members are
the members of \(\N\), together with \(\N\) and \(\N+1\).

\emph{\(G\) is a functional rule for \(\delta\).} For every function \(g\),
the set of values \(\{g(\beta):\beta\in\dom g\}\) exists by Replacement on
\(\dom g\), its union exists by Union, the intersection with \(\N\) by
Separation, and \(\Phi\) of a subset of \(\N\) by Pairing, Union and
Replacement. Each step yields a unique set, so \(G(g)\) is uniquely
determined. Let \(F\) be the function given by
\cref{thm:p01-w02-p01-013-recursion}. For \(\alpha<\delta\) write
\(U_\alpha=\bigcup\{F(\beta):\beta<\alpha\}\); by item~1 of
\cref{thm:p01-w02-p01-013-agreement} the values of
\(F\mathbin{\upharpoonright}\alpha\) are the \(F(\beta)\) with
\(\beta<\alpha\), so \(F(\alpha)=\Phi(\N\cap U_\alpha)\).

\emph{An identity.} For every \(n\in\N\),
\(\{0\}\cup\{k+1:k\in n\}=n+1\). For \(n=0\) both sides are \(\{0\}\). If it
holds for \(n\), then
\[
  \begin{aligned}
  \{0\}\cup\{k+1:k\in n+1\}
  &=\bigl(\{0\}\cup\{k+1:k\in n\}\bigr)\cup\{n+1\}\\
  &=(n+1)\cup\{n+1\}=(n+1)+1,
  \end{aligned}
\]
because the members of \(n+1\) are the members of \(n\) and \(n\) itself.
Hence \(\Phi(n)=n+1\) for \(n\in\N\), using \(n\subseteq\N\).

\emph{Finite stages.} Let \(B\) be the set of \(n\in\N\) with
\(F(n)\ne n+1\), formed by Separation, and suppose \(B\ne\varnothing\). Let
\(n\) be its \(\in\)-least member in the well-ordering of \(\N\). Then
\(F(k)=k+1\) for every \(k\in n\), since such \(k\) lie in \(\N\) and not in
\(B\). We show \(U_n=n\). If \(n=0\), then \(U_0=\bigcup\varnothing=0\).
Otherwise \(n=m+1\) with \(m\in\N\), and \(U_n=\bigcup\{k+1:k\in m+1\}\). It
contains \(m+1\), the term for \(k=m\). Each term is a subset of \(m+1\):
for \(k\in m+1\), either \(k=m\), or \(k\in m\) and then \(k\subseteq m\)
because \(m\) is transitive; in both cases \(k\subseteq m+1\) and
\(k\in m+1\), so \(k+1=k\cup\{k\}\subseteq m+1\). Hence \(U_n=m+1=n\). Since
\(n\subseteq\N\), \(F(n)=\Phi(\N\cap n)=\Phi(n)=n+1\), contradicting
\(n\in B\). So \(F(n)=n+1\) for every \(n\in\N\).

\emph{Stage \(\N\).} \(U_\N=\bigcup\{n+1:n\in\N\}\). Each \(n+1\) is a
member of \(\N\), hence a subset of \(\N\) by transitivity, so
\(U_\N\subseteq\N\); and each \(n\in\N\) lies in \(n+1\), so
\(\N\subseteq U_\N\). Thus \(F(\N)=\Phi(\N)=\{0\}\cup\{k+1:k\in\N\}\). Every
member of this set lies in \(\N\), and every \(n\in\N\) is \(0\) or \(m+1\)
with \(m\in\N\), so it lies in the set. Hence \(F(\N)=\N\).

\emph{Stage \(\N+1\).} The members of \(\N+1\) are the \(n\in\N\) and \(\N\),
so \(U_{\N+1}=U_\N\cup F(\N)=\N\cup\N=\N\), and
\(F(\N+1)=\Phi(\N)=\N\).

\emph{Closure.} For \(n\in\N\), \(\Phi(F(n))=\Phi(n+1)=(n+1)+1\). This set
contains \(n+1\), while \(n+1\notin n+1\) by item~1; so
\(\Phi(F(n))\ne F(n)\). Finally \(\Phi(F(\N))=\Phi(\N)=\N=F(\N)\), as
computed above.
\end{proof}
\end{annexitem}

\begin{annexitem}{annex:p01-w02-p01-013-ill-founded}{thm:p01-w02-p01-013-ill-founded}{Ill-founded recursion}
\begin{proof}
Here \(0=\varnothing\) and \(1=\{0\}\), so \(0\ne1\). For each \(n\in\N\)
the set \(T_n\) is nonempty, because \(n\in n+1\) and \(n+1\in\N\).
For a function \(F\) with domain \(\N\), the argument of item~1 of
\cref{thm:p01-w02-p01-013-agreement} shows that
\(F\mathbin{\upharpoonright}T_n\) is a function with domain \(T_n\) that
agrees with \(F\) there, so its set of values is \(\{F(m):m\in T_n\}\).

\emph{Item~1.} Let \(c\in\{0,1\}\) and let \(F\) be the constant function
with domain \(\N\) and value \(c\). Since \(T_n\ne\varnothing\), the set of
values of \(F\mathbin{\upharpoonright}T_n\) is \(\{c\}\). If \(c=1\), it
contains \(1\) and \(G_1\) gives \(1=F(n)\). If \(c=0\), it does not contain
\(1\), and \(G_1\) gives \(0=F(n)\). So both constant functions satisfy the
equation at every \(n\), and they are different because \(0\ne1\).

\emph{Item~2.} Suppose \(F\) has domain \(\N\) and
\(F(n)=G_2(F\mathbin{\upharpoonright}T_n)\) for every \(n\). Every value of
\(G_2\) is \(0\) or \(1\), so every value of \(F\) is \(0\) or \(1\).

First suppose \(F(n)=1\) for some \(n\). Then \(G_2\) gives \(1\) at
\(F\mathbin{\upharpoonright}T_n\), so \(F(m)\ne1\) for every \(m\in T_n\).
In particular \(n+1\in T_n\), so \(F(n+1)=0\). Then \(G_2\) gives \(0\) at
\(F\mathbin{\upharpoonright}T_{n+1}\), so \(F(m)=1\) for some \(m\in T_{n+1}\).
For that \(m\) we have \(n+1\in m\), hence \(n\in m\) because \(m\) is
transitive and \(n\in n+1\); so \(m\in T_n\) and \(F(m)=1\), a
contradiction.

So \(F(n)=0\) for every \(n\in\N\). Then the set of values of
\(F\mathbin{\upharpoonright}T_0\) is \(\{0\}\), which does not contain
\(1\), and \(G_2\) gives \(F(0)=1\), a contradiction. Hence no such \(F\)
exists.
\end{proof}
\end{annexitem}
